% Generated draft: requires mathematical and source review.
\documentclass[11pt]{article}
\usepackage[margin=25mm]{geometry}
\usepackage{fontspec}
\setmainfont{Latin Modern Roman}
\usepackage{amsmath,amssymb,mathtools,unicode-math}
\setmathfont{Latin Modern Math}
\usepackage{enumitem,needspace,xurl,hyperref}
\setlength{\parindent}{0pt}
\setlength{\parskip}{5pt}
\setlength{\emergencystretch}{3em}
\newcommand{\N}{\mathbb{N}}
\newcommand{\Z}{\mathbb{Z}}
\newcommand{\Q}{\mathbb{Q}}
\newcommand{\R}{\mathbb{R}}
\newcommand{\C}{\mathbb{C}}
\newcommand{\F}{\mathbb{F}}
\newcommand{\A}{\mathbb{A}}
\newcommand{\PP}{\mathbb{P}}
\newcommand{\E}{\mathbb{E}}
\newcommand{\eps}{\varepsilon}
\newcommand{\dd}{\,\mathrm{d}}
\begin{document}
\section*{UnsolvedMath ID: 11000109}
\textbf{Source problem code:} catalogue entry
\subsection{Definitions and mathematical statement}
Let $G=\operatorname{Mod}_{1,2}$. In the customary finite-type surface convention this is the mapping class group of a genus-one surface with two punctures or marked points; the excerpt does not specify whether they may be permuted, whether orientation reversal is allowed, or whether a boundary convention is intended instead.

\textbf{Problem 2.4.} For every subgroup $\Gamma\le G$ with $[G:\Gamma]<\infty$ and every group automorphism $\varphi:\Gamma\to\Gamma$, does there exist $\Phi\in\operatorname{Aut}(G)$ such that
\[
 \Phi|_{\Gamma}=\varphi?
\]
In particular, an extension must preserve $\Gamma$ setwise. Equivalently, is the restriction map
\[
 \{\Phi\in\operatorname{Aut}(G):\Phi(\Gamma)=\Gamma\}
 \longrightarrow\operatorname{Aut}(\Gamma)
\]
surjective for every finite-index $\Gamma$? No normality or torsion-freeness hypothesis is imposed on $\Gamma$. This is an extension to an abstract automorphism of the full group, not an assertion that every automorphism of $\Gamma$ must be inner in $\Gamma$.

As subsequent context, the source reports a theorem of Harvey and the chapter author [14]: for $g>h$, a homomorphism $\operatorname{Mod}_g\to\operatorname{Mod}_h$ has trivial image except that, for $g=2$, its image is allowed to have order at most two. That reported result concerns homomorphisms between different genus groups; it is not the finite-index extension assertion posed here. The source then points to a following question, whose text is not included in this record.
\subsection{Short English statement}
Does every automorphism of every finite-index subgroup of $\operatorname{Mod}_{1,2}$ extend to an automorphism of the whole group? The exact puncture and orientation conventions for this exceptional surface group need confirmation.
\subsection{Sources}
\begin{itemize}
\item Local UnsolvedMath registry, permanent ID 11000109, source code catalogue entry; no individual proposer or proposal year supplied.
\item Farb, edited volume \emph{Problems on Mapping Class Groups and Related Topics}, Problem 2.4, PDF page 94; \url{https://www.math.uchicago.edu/~farb/papers/mcgbook.pdf}. The registry records PDF-font-structure extraction and access on 2026-07-29; its inventory distinguishes 329 labeled headings by source page.
\item Harvey and the chapter author, source reference [14], for the subsequently reported different-genus homomorphism theorem. The full citation and exact mapping-class conventions are not supplied.
\end{itemize}
\end{document}
